% =============================================================================
%  Capítulo 4: Regresión por Mínimos Cuadrados
%  Basado en:
%    Steven L. Brunton, "Optimization: A Bootcamp for Machine Learning,
%    Inverse Problems, and Control", Cambridge University Press, 2026.
%    Capítulo 4, secciones 4.1--4.8 (pp. 122-166).
%
%  Contenido verificado del capítulo:
%    4.1 Least-Squares Regression
%    4.2 Singular Value Decomposition
%    4.3 Minimum Norm Regression
%    4.4 Condition Number and Regularization
%    4.5 Constrained Least-Squares
%    4.6 Quadratic Programming
%    4.7 Nonlinear Least-Squares: Gauss-Newton Methods
%    4.8 The Conjugate Gradient Method
%
%  Compilar:  make
% =============================================================================

\documentclass[10pt,aspectratio=169]{beamer}

% --- Idioma y codificación ---------------------------------------------------
\usepackage[utf8]{inputenc}
\usepackage[T1]{fontenc}
\usepackage[spanish]{babel}

% --- Matemáticas -------------------------------------------------------------
\usepackage{amsmath,amssymb,mathtools}

% --- Tablas y diseño ---------------------------------------------------------
\usepackage{booktabs}
\usepackage{tabularx}
\usepackage{array}
\usepackage{lmodern}
\usepackage{microtype}
\usepackage{graphicx}

% --- Colores -----------------------------------------------------------------
\definecolor{unah}{RGB}{0,82,155}
\definecolor{acento}{RGB}{0,135,110}

\setbeamercolor{structure}{fg=unah}
\setbeamercolor{frametitle}{fg=unah}
\setbeamercolor{itemize item}{fg=acento}
\setbeamercolor{itemize subitem}{fg=unah}
\setbeamertemplate{itemize item}{\color{acento}\scriptsize$\blacktriangleright$}
\setbeamertemplate{itemize subitem}{\color{unah}\scriptsize$\bullet$}

% --- Plantillas --------------------------------------------------------------
\setbeamertemplate{navigation symbols}{}
\setbeamertemplate{footline}[frame number]

% --- Atajos ------------------------------------------------------------------
\newcommand{\fuente}{\vfill{\scriptsize\color{gray}Brunton, \emph{Optimization Bootcamp}, Cap.~4}}
\newcommand{\R}{\mathbb{R}}
\newcommand{\col}{\mathrm{col}}
\newcommand{\rk}{\mathrm{rank}}
\newcommand{\argmin}{\operatorname*{arg\,min}}

\title{Cap. 4: Regresión por Mínimos Cuadrados}
\subtitle{Optimización Numérica --- Maestría en Matemática}
\author{Luis Berlioz}
\date{\today}

% =============================================================================
\begin{document}
% =============================================================================

% --- Portada -----------------------------------------------------------------
\begin{frame}
  \titlepage
\end{frame}

% --- Esquema -----------------------------------------------------------------
\begin{frame}{Temas del capítulo}
  \tableofcontents
\end{frame}

% =============================================================================
\section{Mínimos cuadrados lineales}
% =============================================================================

\begin{frame}{El problema en una línea}
  \small
  Un sistema lineal \textbf{con más ecuaciones que incógnitas} no tiene solución
  exacta. Se convierte en un problema de optimización:

  \[
    Ax = b \qquad\text{\scriptsize\textcolor{gray}{(4.1)}}
    \qquad\Longrightarrow\qquad
    \min_x \; \|Ax - b\|_2^2 \qquad\text{\scriptsize\textcolor{gray}{(4.2)}}
  \]

  \begin{itemize}
    \item La solución es \textbf{analítica y cerrada}:
      \[
        x = A^{\dagger} b \tag{4.3}
      \]
      con $A^{\dagger}$ la \textbf{pseudo-inversa de Moore--Penrose}.
    \item La generalización no lineal
      $\displaystyle \min_\theta \|N_\theta(x) - y\|_2^2$
      \textbf{no es convexa} y exige métodos iterativos ($\S4{.}7$).
  \end{itemize}

  \begin{block}{La cadena del capítulo}
    \centering\scriptsize
    solución cerrada $\to$ SVD $\to$ norma mínima $\to$ conditioning
    $\to$ regularización $\to$ restricciones y QP $\to$ Gauss--Newton
    $\to$ gradiente conjugado
  \end{block}
\end{frame}

\begin{frame}{Cuadrado, sobredeterminado, subdeterminado}
  \begin{columns}[T,onlytextwidth]
    \begin{column}{.40\textwidth}
      \small
      \begin{itemize}
        \item \textbf{Cuadrado} ($m=n$): invertir la matriz; solución exacta.
        \item \textbf{Sobredeterminado} ($m \gg n$): más ecuaciones que
          incógnitas. \emph{Tall-skinny}. Casi nunca hay solución exacta
          $\Rightarrow$ se minimiza el error.
        \item \textbf{Subdeterminado} ($m \ll n$): menos ecuaciones que
          incógnitas. \emph{Short-fat}. Infinitas soluciones
          $\Rightarrow$ se elige la de norma mínima.
      \end{itemize}
    \end{column}
    \begin{column}{.58\textwidth}
      \begin{figure}
        \centering
        \includegraphics[width=\textwidth]{images/fig_4_2.png}
      \end{figure}
      \vspace{-2mm}
      {\scriptsize\color{gray}Fig.~4.2: regresión sobredeterminada (izq.)
      y subdeterminada (der.)}
    \end{column}
  \end{columns}
  \fuente
\end{frame}

\begin{frame}{El modelo de regresión lineal}
  \small
  Los datos se ordenan como filas de una matriz $A \in \R^{m\times n}$ y un
  vector de resultados $b \in \R^m$. Se busca el peso $x \in \R^n$ tal que
  $b$ se aproxime por combinación lineal de las \textbf{columnas} de $A$:

  \[
    Ax \approx b
  \]

  \begin{itemize}
    \item Cada columna de $A$ es un \textbf{atributo} o predictor; cada fila,
      una \textbf{medición}.
    \item $b$ casi nunca es exacto: viene contaminado con ruido de medición,
      modelado como ruido blanco Gaussiano i.i.d.
      \[
        b = b_{\text{true}} + \epsilon
      \]
    \item Motivación física (p.~122): el equilibrio de fuerzas en las juntas de
      un puente da una ecuación por junta. Con ruido en $b$ o imperfecciones en
      $A$, resolver por las fuerzas internas es un \textbf{problema inverso de
      mínimos cuadrados}.
  \end{itemize}
  \fuente
\end{frame}

\begin{frame}{Las ecuaciones normales}
  \small
  Expandir el objetivo y anular el gradiente:

  \begin{align*}
    \|Ax-b\|_2^2 &= (Ax-b)^T(Ax-b) \\
                 &= x^TA^TAx - x^TA^Tb - b^TAx + b^Tb \tag{4.5c}
  \end{align*}
  \[
    \nabla_x \|Ax-b\|_2^2 = 2A^TAx - 2A^Tb = 0
    \;\;\Longrightarrow\;\;
    \boxed{\,A^TAx = A^Tb\,} \tag{4.6b}
  \]
  \[
    x = (A^TA)^{-1}A^Tb \tag{4.6c}
  \]

  \begin{alertblock}{La trampa}
    (4.6c) \textbf{solo vale si $A^TA$ es invertible}, es decir si $A$ tiene
    $n$ columnas linealmente independientes y $m \ge n$. Es la fórmula del caso
    \textbf{sobredeterminado}; no sirve para \emph{short-fat}.
  \end{alertblock}
  \fuente
\end{frame}

\begin{frame}{La pseudo-inversa de Moore--Penrose}
  \small
  \[
    A^{\dagger} = (A^TA)^{-1}A^T
    \qquad\text{(izquierda)}
    \qquad\qquad
    A^{\dagger} = A^T(AA^T)^{-1}
    \qquad\text{(derecha)}
  \]

  \begin{itemize}
    \item Está definida incluso para $A$ no cuadrada, pero cada forma exige
      rango completo en su lado.
    \item Numéricamente \textbf{no se invierte $A^TA$}: su número de condición
      es el cuadrado del de $A$. Se usa la \textbf{SVD}:
      \[
        A = U\Sigma V^T
        \quad\Longrightarrow\quad
        A^{\dagger} = V\Sigma^{-1}U^T
        \quad\Longrightarrow\quad
        x = V\Sigma^{-1}U^T b \tag{4.7}
      \]
    \item $U$ y $V$ son ortogonales: invertirlas es transponerlas.
      $\Sigma$ es diagonal: invertirla son $n$ divisiones.
  \end{itemize}

  \begin{block}{Por qué importa}
    \scriptsize
    La SVD es la técnica de elección para invertir matrices grandes: es
    numéricamente estable, siempre existe y no necesita formar $A^TA$.
  \end{block}
  \fuente
\end{frame}

\begin{frame}{Norma 2 frente a norma 2 al cuadrado}
  \begin{columns}[c,onlytextwidth]
    \begin{column}{.52\textwidth}
      \begin{figure}
        \centering
        \includegraphics[width=\textwidth]{images/fig_4_3.png}
      \end{figure}
    \end{column}
    \begin{column}{.45\textwidth}
      \small
      \begin{itemize}
        \item Izquierda: $\|x\|_2 = \sqrt{x_1^2+x_2^2}$.
          Sus conjuntos de nivel son \textbf{círculos}.
        \item Derecha: $\|x\|_2^2 = x_1^2+x_2^2$.
          Mismos conjuntos de nivel, misma solución óptima.
        \item Minimizar $\|\cdot\|_2$ o $\|\cdot\|_2^2$ da \textbf{el mismo
          $x^*$}: el cuadrado es monótono.
        \item Pero el cuadrado \textbf{pesa distinto} los errores: un error
          grande cuesta mucho más que dos pequeños. De ahí la sensibilidad a
          outliers.
      \end{itemize}
    \end{column}
  \end{columns}
  \fuente
\end{frame}

\begin{frame}[fragile]{Ejemplo: datos aleatorios (E4.1)}
  \small
  Pseudo-inversa calculada por SVD frente a \texttt{np.linalg.pinv}:

  \begin{semiverbatim}
import numpy as np
np.random.seed(42)
m, n = 20, 5
A = np.random.randn(m, n)
b = np.random.randn(m)

U, S, VT = np.linalg.svd(A, full_matrices=0)
x = VT.T @ np.linalg.inv(np.diag(S)) @ U.T @ b

x = np.linalg.pinv(A) @ b          # alternativa
  \end{semiverbatim}

  \begin{itemize}
      \item \texttt{full\_matrices=0} encuentra la descomposición \textbf{SVD económica}: solo los
      valores singulares no nulos, matrices $m\times n$ en vez de $m\times m$.
    \item Ambas líneas dan el mismo $x$ hasta error de redondeo.
  \end{itemize}
  \fuente
\end{frame}

\begin{frame}{Ejemplo: calor del cemento (E4.2)}
  \begin{columns}[T,onlytextwidth]
    \begin{column}{.52\textwidth}
      \begin{figure}
        \centering
        \includegraphics[width=\textwidth]{images/fig_4_5.png}
      \end{figure}
      \vspace{-2mm}
      {\scriptsize\color{gray}Fig.~4.5: calor generado por 13 mezclas de cemento.}
    \end{column}
    \begin{column}{.45\textwidth}
      \small
      $A \in \R^{13\times 4}$: cuatro ingredientes, trece mezclas. Se buscan los
      pesos $x$ que relacionan proporciones con calor.

      \vspace{2mm}
      \centering
      \begin{tabular}{@{}lr@{}}
        \toprule
        Ingrediente & $x_j$ \\
        \midrule
        Tricalcium aluminate   & 2.1930 \\
        Tricalcium silicate    & 1.1533 \\
        Tetracalcium alumifer. & 0.7585 \\
        Beta-dicalcium silicate& 0.4863 \\
        \bottomrule
      \end{tabular}

      \vspace{3mm}
      \raggedright
      \scriptsize
      Es el experimento clásico de Hald (1952): cuatro predictores fuertemente
      correlacionados y solo trece datos.
    \end{column}
  \end{columns}
  \fuente
\end{frame}

\begin{frame}{Ejemplo: precios de vivienda en Boston (E4.3)}
  \small
  \begin{columns}[T,onlytextwidth]
    \begin{column}{.58\textwidth}
      \begin{itemize}
        \item $506$ viviendas, $13$ atributos: criminalidad per cápita, tasa de
          impuesto predial, etc.
        \item $A \in \R^{506\times 13}$: problema claramente sobredeterminado.
        \item \textbf{Hay que añadir una columna de unos}: da el término
          constante (la ordenada al origen). Equivale a centrar cada columna
          restando su media.
        \item El ajuste no captura los precios extremos altos: los
          \textbf{modelos lineales simples fallan en las colas}.
      \end{itemize}
    \end{column}
    \begin{column}{.40\textwidth}
      \begin{figure}
        \centering
        \includegraphics[width=\textwidth]{images/fig_4_6.png}
      \end{figure}
      \vspace{-2mm}
      {\scriptsize\color{gray}Fig.~4.6: valor real frente a valor ajustado.}
    \end{column}
  \end{columns}
  \fuente
\end{frame}

\begin{frame}{Outliers: por qué $\ell_2$ no basta}
  \begin{columns}[T,onlytextwidth]
    \begin{column}{.50\textwidth}
      \small
      \begin{itemize}
        \item Un outlier entra \textbf{al cuadrado}: su error domina y
          \emph{arrastra} toda la recta.
        \item Alternativa robusta: \textbf{desviación absoluta mínima}
          \[
            \min_x \; \|Ax - b\|_1 \tag{4.8}
          \]
        \item Regla general: usar $\ell_1$ como pérdida suele ser más robusto
          a outliers. El precio es perder la solución cerrada: hay que resolver
          un problema de optimización.
      \end{itemize}
    \end{column}
    \begin{column}{.48\textwidth}
      \begin{figure}
        \centering
        \includegraphics[width=\textwidth]{images/fig_4_7.png}
      \end{figure}
      \vspace{-2mm}
      {\scriptsize\color{gray}Fig.~4.7: $\ell_2$ (sensible) frente a
      $\ell_1$ (robusta).}
    \end{column}
  \end{columns}
  \fuente
\end{frame}

\begin{frame}{Panel 4.1: ¿mínimos cuadrados o desviación absoluta?}
  \small
  Doscientos años de debate, con una lectura estadística limpia:

  \vspace{2mm}
  \centering
  \begin{tabular}{@{}lll@{}}
    \toprule
    & \textbf{Mínimos cuadrados} ($\ell_2$) & \textbf{LAD} ($\ell_1$) \\
    \midrule
    Estima               & la \textbf{media}   & la \textbf{mediana} \\
    Robustez a outliers  & baja                & alta \\
    Óptimo               & cerrado (ecuaciones normales) & programación lineal \\
    Historia             & Gauss y Legendre, $\sim$1800 & Boscovich, 1757 \\
    \bottomrule
  \end{tabular}

  \vspace{4mm}
  \raggedright
  \begin{itemize}
    \item La media es insesgada para ruido Gaussiano; la mediana es insesgada
      para ruido Laplaciano.
    \item Ejemplo del libro (sep.\ 2024): riqueza mediana en EE.\,UU.\
      $\sim\$200$k frente a riqueza media $\sim\$1{,}000$k.
      \textbf{La media la mueven los outliers; la mediana no.}
  \end{itemize}
  \fuente
\end{frame}

\begin{frame}{Los cuatro subespacios fundamentales}
  \small
  \begin{columns}[T,onlytextwidth]
    \begin{column}{.48\textwidth}
      \begin{itemize}
        \item $\col(A)$: span de las columnas (\emph{rango}).
        \item $\mathrm{row}(A) = \col(A^T)$: span de las filas.
        \item $\ker(A)$: vectores con $Ax=0$.
        \item $\ker(A^T)$: complemento ortogonal de $\col(A)$.
      \end{itemize}
    \end{column}
    \begin{column}{.48\textwidth}
      \begin{block}{Descomposición en suma directa (p.~131)}
        \[
          \col(A) \oplus \ker(A^T) = \R^{m} \tag{4.9a}
        \]
        \[
          \col(A^T) \oplus \ker(A) = \R^{n} \tag{4.9b}
        \]
      \end{block}
    \end{column}
  \end{columns}

  \vspace{2mm}
  Consecuencias sobre $Ax=b$:
  \begin{itemize}
    \item $b \in \col(A)$ y $\dim\ker(A) \ne 0$ $\Rightarrow$ \textbf{infinitas
      soluciones} (garantizado si $A$ es short-fat).
    \item $b \notin \col(A)$ $\Rightarrow$ \textbf{ninguna solución}: el sistema
      es inconsistente.
    \item Se elige entonces \textbf{la mejor aproximación} (\emph{least squares})
      o \textbf{la de norma mínima}, según el régimen.
  \end{itemize}
  \fuente
\end{frame}

\begin{frame}{Las dos soluciones canónicas}
  \small
  \begin{columns}[T,onlytextwidth]
    \begin{column}{.48\textwidth}
      \begin{block}{Sobredeterminado: mínimos cuadrados}
        \[
          \min_x \; \|Ax-b\|_2^2
        \]
        \[
          x = A^{\dagger}b = (A^TA)^{-1}A^Tb
        \]
        \scriptsize Minimiza el error; $b$ se proyecta sobre $\col(A)$.
      \end{block}
    \end{column}
    \begin{column}{.48\textwidth}
      \begin{block}{Subdeterminado: norma mínima}
        \[
          \min_x \; \|x\|_2^2 \quad \text{s.a. } Ax=b
        \]
        \[
          x = A^T(AA^T)^{-1}b
        \]
        \scriptsize Satisface la ecuación exacta con el $x$ más corto posible.
      \end{block}
    \end{column}
  \end{columns}

  \vspace{1mm}
  \begin{itemize}
    \item Ambas son la \textbf{misma} pseudo-inversa $A^{\dagger}$, aplicada por
      lados distintos.
    \item Obs.~1 (p.~131): con matrices \textbf{aleatorias} los estereotipos son
      ciertos con altísima probabilidad. Para $A$ aleatoria con $m \gg n$ es
      casi imposible que $b \in \col(A)$.
  \end{itemize}
  \fuente
\end{frame}

% =============================================================================
\section{La descomposición en valores singulares}
% =============================================================================

\begin{frame}{Qué es la SVD}
  \small
  \[
    A = U\Sigma V^T \tag{4.10}
  \]
  \begin{columns}[T,onlytextwidth]
    \begin{column}{.50\textwidth}
      \begin{itemize}
        \item $U \in \R^{m\times m}$ y $V \in \R^{n\times n}$
          \textbf{ortogonales}: columnas ortonormales.
        \item Columnas de $U$: \textbf{vectores singulares izquierdos}.
        \item Columnas de $V$: \textbf{derechos}.
        \item Diagonal de $\Sigma$: valores singulares ordenados
          \[
            \sigma_1 \ge \sigma_2 \ge \cdots \ge \sigma_n \ge 0 .
          \]
        \item El \textbf{rango} de $A$ es el número de $\sigma_j$ no nulos.
        \item Para matrices complejas, $U,V$ son unitarias y $T$ pasa a ser
          transpuesta conjugada $*$.
      \end{itemize}
    \end{column}
    \begin{column}{.48\textwidth}
      \begin{figure}
        \centering
        \includegraphics[width=\textwidth]{images/fig_4_8.png}
      \end{figure}
      \vspace{-2mm}
      {\scriptsize\color{gray}Fig.~4.8: la SVD descompone en filas y columnas
      ortogonales ordenadas por varianza.}
    \end{column}
  \end{columns}
  \fuente
\end{frame}

\begin{frame}{SVD económica y truncada}
  \small
  \begin{columns}[T,onlytextwidth]
    \begin{column}{.52\textwidth}
      \textbf{SVD económica} (solo $\sigma$ no nulos, $m \ge n$):
      \[
        A = \hat{U}\hat{\Sigma}V^T \tag{4.11}
      \]
      \textbf{SVD truncada de rango $r$}:
      \[
        \tilde{A} = \tilde{U}\tilde{\Sigma}\tilde{V}^T
      \]
      $\tilde{U},\tilde{V}$: primeras $r$ columnas;
      $\tilde{\Sigma}$: bloque $r\times r$.

      \vspace{2mm}
      \begin{itemize}
        \item ``Truncar'' = quedarse con los $r$ patrones dominantes.
        \item Es el mismo mecanismo que \textbf{PCA}.
      \end{itemize}
    \end{column}
    \begin{column}{.46\textwidth}
      \begin{figure}
        \centering
        \includegraphics[width=\textwidth]{images/fig_4_9.png}
      \end{figure}
      \vspace{-2mm}
      {\scriptsize\color{gray}Fig.~4.9: aproximación truncada con
      $r$ valores singulares.}
    \end{column}
  \end{columns}
  \fuente
\end{frame}

\begin{frame}{La SVD como suma de matrices de rango 1}
  \small
  \begin{block}{Panel 4.2 (p.~133)}
    Como $\Sigma$ es diagonal, la SVD se lee como una suma de términos
    $\sigma_j u_j v_j^T$, cada uno de rango 1 y ortogonal a los demás:
    \[
      X \approx \sigma_1 u_1 v_1^T + \sigma_2 u_2 v_2^T
             + \cdots + \sigma_r u_r v_r^T
    \]
    Ordenados por magnitud, los primeros términos capturan lo grueso;
    los últimos, el detalle fino (y el ruido).
  \end{block}

  \vspace{1mm}
  \begin{columns}[T,onlytextwidth]
    \begin{column}{.48\textwidth}
      \textbf{Subespacios en términos de $U,V$:}
      \begin{itemize}
        \item $\col(A) = \col(\hat{U})$
        \item Complemento ortogonal de $\col(A)$: $\col(\hat{U}^\perp)$
        \item $\mathrm{row}(A)$ generado por columnas de $V$
        \item $\ker(A) = \col(V^\perp)$
      \end{itemize}
    \end{column}
    \begin{column}{.48\textwidth}
      \begin{block}{Por qué es la herramienta}
        \scriptsize
        Da a la vez: rango numérico, base de los cuatro subespacios,
        pseudo-inversa estable, aproximación óptima de bajo rango y
        diagnóstico de conditioning. Todo con una sola factorización.
      \end{block}
    \end{column}
  \end{columns}
  \fuente
\end{frame}

\begin{frame}{Resolver $Ax=b$ con la SVD}
  \small
  Sustituyendo la SVD truncada exacta ($A=\hat{U}\hat{\Sigma}V^T$) y
  ``invirtiendo'' cada factor:

  \[
    A^{\dagger} \triangleq V\hat{\Sigma}^{-1}\hat{U}^T
    \qquad\Longrightarrow\qquad
    A^{\dagger}A = VV^T \tag{4.12}
  \]
  \[
    A^{\dagger}Ax^* = A^{\dagger}b
    \;\;\Longrightarrow\;\;
    x^* = V\hat{\Sigma}^{-1}\hat{U}^T b \tag{4.13}
  \]

  \begin{itemize}
    \item Sustituyendo de vuelta:
      $Ax^* = \hat{U}\hat{U}^T b$ \ \ (4.14b).
    \item \textbf{Interpretación: $x^*$ es la proyección de $b$ sobre
      $\col(\hat{U})$}. El residuo $b - \hat{U}\hat{U}^Tb$ es la parte de $b$
      que vive en $\ker(A^T)$ y que \emph{ningún} $x$ puede explicar.
    \item $A^{\dagger}A = I$ solo si la SVD captura \emph{todos} los valores
      singulares no nulos. Con rango $r$ se obtiene $\tilde{V}\tilde{V}^T$:
      una aproximación de la identidad.
  \end{itemize}
  \fuente
\end{frame}

\begin{frame}{Teorema de Eckart--Young}
  \small
  \begin{block}{Teorema 1 (Eckart--Young [159], p.~135)}
    La mejor aproximación de rango $r$ en norma de Frobenius es la SVD truncada:
    \[
      \argmin_{\tilde{X}:\ \rk(\tilde{X})=r} \|X - \tilde{X}\|_F
      = \tilde{U}\tilde{\Sigma}\tilde{V}^T \tag{4.15}
    \]
    con $\|X\|_F = \big(\sum_{ij} X_{ij}^2\big)^{1/2}$, equivalente a la
    norma 2 de $X(:)$.
  \end{block}

  \begin{columns}[T,onlytextwidth]
    \begin{column}{.48\textwidth}
      \textbf{Error exacto en Frobenius:}
      \[
        \|X - \tilde{X}\|_F^2 = \sum_{k=r+1}^{m} \sigma_k^2 \tag{4.16}
      \]
      \scriptsize Toda matriz de rango $r$ tiene al menos este error.
      El error relativo (4.17) mide la \textbf{varianza faltante} --- la
      métrica que usa PCA.
    \end{column}
    \begin{column}{.48\textwidth}
      \textbf{Sorprendentemente, también en norma espectral:}
      \[
        \|X - \tilde{X}\|_2 = \sigma_{r+1} \tag{4.19}
      \]
      \scriptsize porque
      $X-\tilde{X} = \sum_{k>r}\sigma_k u_k v_k^T$ y el máximo se alcanza en
      $v = v_{r+1}$.
    \end{column}
  \end{columns}

  \vspace{1mm}
  \footnotesize
  Schmidt (Gram--Schmidt) ya lo había probado para espacios de funciones;
  Eckart y Young lo redescubrieron en 1936.
  \fuente
\end{frame}

\begin{frame}{Costo de la SVD}
  \small
  \begin{columns}[T,onlytextwidth]
    \begin{column}{.48\textwidth}
      \textbf{Algoritmo (p.~134):}
      \begin{enumerate}
        \item Reducir $A$ a \textbf{bidiagonal}. Con $m \gg n$: primero
          factorización \textbf{QR}, luego reflejos de Householder.
        \item SVD de la bidiagonal con el \textbf{QR modificado de
          Golub--Kahan}.
      \end{enumerate}
    \end{column}
    \begin{column}{.48\textwidth}
      \textbf{Entonces $A^{\dagger}$ es barata:}
      \begin{itemize}
        \item invertir $U$ o $V$: multiplicar por la transpuesta,
          $O(m^2)$;
        \item invertir $\Sigma$: $O(m)$, es diagonal;
        \item \emph{contraste:} invertir una matriz densa cuesta $O(m^3)$.
      \end{itemize}
    \end{column}
  \end{columns}

  \vspace{2mm}
  \begin{block}{Regla práctica}
    \scriptsize
    Nunca se forma ni se invierte $A^TA$: se factoriza $A$ una vez y se
    reutiliza la misma SVD para mínimos cuadrados, norma mínima,
    regularización y truncamiento.
  \end{block}
  \fuente
\end{frame}

% =============================================================================
\section{Norma mínima y compressed sensing}
% =============================================================================

\begin{frame}{El problema de norma mínima}
  \small
  Caso \textbf{subdeterminado}: hay infinitas soluciones de $Ax=b$.
  Se elige la más corta:

  \[
    \min_x \; \|x\|_2^2 \qquad \text{s.a. } Ax = b \tag{4.22a--b}
  \]

  \begin{columns}[T,onlytextwidth]
    \begin{column}{.48\textwidth}
      \begin{itemize}
        \item Aquí el objetivo es $\|x\|_2^2$, \textbf{no} el residual.
          La ecuación entra como \textbf{restricción}.
        \item Geométricamente: la bola $\ell_2$ de radio mínimo que toca el
          hiperplano $Ax=b$.
        \item La solución $\ell_2$ mínima es \textbf{densa}: reparte el peso
          entre todas las coordenadas.
      \end{itemize}
    \end{column}
    \begin{column}{.48\textwidth}
      \begin{figure}
        \centering
        \includegraphics[width=\textwidth]{images/fig_4_11.png}
      \end{figure}
      \vspace{-2mm}
      {\scriptsize\color{gray}Fig.~4.11: las bolas $\ell_2$, $\ell_1$ y
      $\ell_\infty$ cortando la recta $Ax=b$.}
    \end{column}
  \end{columns}
  \fuente
\end{frame}

\begin{frame}{Derivación por multiplicadores de Lagrange}
  \small
  \[
    L(x,\lambda) = x^Tx + \lambda^T(Ax-b) \tag{4.23}
  \]
  \[
    \nabla_x L = 2x + A^T\lambda = 0
    \;\;\Longrightarrow\;\;
    x = -\tfrac12 A^T\lambda \tag{4.24}
  \]
  \[
    \nabla_\lambda L = Ax - b = 0
    \;\;\Longrightarrow\;\;
    -AA^T\lambda = b
    \;\;\Longrightarrow\;\;
    \lambda = -2(AA^T)^{-1}b \tag{4.25}
  \]
  \[
    \boxed{\,x = A^T(AA^T)^{-1}b\,} \tag{4.26--4.27}
  \]

  \begin{block}{Pseudo-inversa derecha}
    \scriptsize
    Requiere que $AA^T$ sea invertible, es decir que $A$ tenga \textbf{rango de
    fila completo} (lo genérico en una matriz short-fat). Con la SVD es la misma
    fórmula $x = V\Sigma^{-1}U^Tb$: $A^{\dagger}$ es \emph{una sola} matriz.
  \end{block}
  \fuente
\end{frame}

\begin{frame}{Compressed sensing: la idea}
  \small
  \begin{itemize}
    \item Imágenes y audio son \textbf{compresibles}: en una base Fourier o
      wavelet, casi todos los coeficientes son despreciables. Es la base de
      JPEG y MP3.
    \item Pero el paradigma clásico \textbf{mide todo y comprime después}.
    \item Invertirlo: medir solo una \textbf{submuestra aleatoria} y reconstruir
      los pocos modos activos. Del millón de píxeles de una imagen,
      potencialmente basta el $10\%$.
  \end{itemize}

  \vspace{1mm}
  \begin{block}{Formulación}
    Señal compresible $x = \Psi s$, con $s$ \textbf{$K$-disperso}
    ($K$ entradas no nulas) y mediciones $y = Cx$:
    \[
      y = Cx = C\Psi s = \Theta s \tag{4.29}
    \]
    con $C \in \R^{p\times n}$ y $p \ll n$: sistema subdeterminado.
  \end{block}
  \fuente
\end{frame}

\begin{frame}{De $\ell_0$ a $\ell_1$: la relajación convexa}
  \small
  La solución más dispersa es la de \textbf{cardinalidad mínima}:
  \[
    \hat{s} = \argmin_s \; \|s\|_0
    \qquad \text{s.a. } y = C\Psi s \tag{4.30}
  \]
  \begin{itemize}
    \item $\|\cdot\|_0$ cuenta entradas no nulas: \textbf{no convexa} y
      combinatoria. Exige verificar todos los $K$-dispersos posibles
      $\Rightarrow$ \textbf{intratable} para $n$ y $K$ moderados.
    \item Bajo condiciones sobre la matriz de medición $C$, se relaja a un
      problema \textbf{convexo}:
      \[
        \hat{s} = \argmin_s \; \|s\|_1
        \qquad \text{s.a. } y = C\Psi s \tag{4.31}
      \]
  \item \textbf{La solución $\ell_1$ mínima es dispersa; la
      $\ell_2$ mínima no lo es.} La bola $\ell_1$ tiene vértices en los ejes,
      y el óptimo casi siempre cae en uno.
    \item Con ruido, $\|C\Psi s - y\|_2 < \epsilon$ (4.32) o la forma
      penalizada $\|C\Psi s-y\|_2 + \lambda\|s\|_1$ (4.33).
  \end{itemize}
  \fuente
\end{frame}

% =============================================================================
\section{Conditioning y regularización}
% =============================================================================

\begin{frame}{Número de condición: el ruido amplificado}
  \small
  \[
    \kappa(A) = \frac{\sigma_{\max}(A)}{\sigma_{\min}(A)} \tag{4.34}
  \]
  Mide la sensibilidad de $A$ frente a errores en la entrada. Si $b$ viene con
  error $\epsilon_b$:
  \[
    A(x+\epsilon_x) = b + \epsilon_b
    \qquad\Longrightarrow\qquad
    \frac{\|x\|}{\|\epsilon_x\|}
    = \frac{\sigma_{\min}}{\sigma_{\max}}\,
      \frac{\|b\|}{\|\epsilon_b\|}
  \]

  \begin{alertblock}{Regla de oro}
    El ruido en $b$ se amplifica por $\kappa(A)$ al invertir. Con columnas casi
    paralelas, $\kappa(A)$ es enorme y $\kappa(A^TA)=\kappa(A)^2$ es peor:
    por eso \textbf{jamás} se resuelve por ecuaciones normales en ese régimen.
  \end{alertblock}

  \vspace{1mm}
  \small
  \textbf{Mitigación por truncamiento} (p.~140): descartar los $\sigma$ bajo
  $\kappa_{\text{obj}}\cdot\sigma_{\max}$ eleva el $\sigma_{\min}$ efectivo;
  el precio es un subespacio $\hat{U}$ más pequeño.
  \fuente
\end{frame}

\begin{frame}{Regularización de Tikhonov (ridge)}
  \small
  Añadir una penalización sobre $\|x\|_2^2$:
  \[
    x = \argmin_{x'} \; \|Ax'-b\|_2^2 + \alpha\|x'\|_2^2 \tag{4.38}
  \]
  \[
    x = (A^TA+\alpha I)^{-1}A^Tb \tag{4.39}
  \]

  \begin{columns}[T,onlytextwidth]
    \begin{column}{.50\textwidth}
      \begin{itemize}
        \item El término $\alpha I$ \textbf{desplaza los valores propios lejos
          de cero}: el número de condición efectivo es
          $(\sigma_{\max}^2+\alpha)/(\sigma_{\min}^2+\alpha)$.
        \item Se puede ver como un multiplicador de Lagrange que impone
          $\|x\|_2^2 = c$: empuja a que las entradas tengan magnitud parecida
          en vez de grandes y opuestas que se cancelan.
        \item ``Ridge'' en estadística; en inversos es ``Tikhonov''.
      \end{itemize}
    \end{column}
    \begin{column}{.48\textwidth}
      \begin{figure}
        \centering
        \includegraphics[width=\textwidth]{images/fig_4_12.png}
      \end{figure}
      \vspace{-2mm}
      {\scriptsize\color{gray}Fig.~4.12: las soluciones ruidosas (elipses
      azules) se arrastran hacia los círculos de $\|x\|_2^2$.}
    \end{column}
  \end{columns}
  \fuente
\end{frame}

\begin{frame}{El compromiso sesgo--varianza}
  \begin{columns}[T,onlytextwidth]
    \begin{column}{.50\textwidth}
      \small
      \begin{itemize}
        \item \textbf{Sin regularizar:} alta varianza. El elipsoide de
          soluciones posibles ante ruido es enorme y muy excéntrico.
        \item \textbf{Con $\alpha$:} la varianza se encoge, pero la
          distribución se \textbf{desplaza}: aparece sesgo.
        \item $\alpha \to 0$: mínimos cuadrados puros.
          $\alpha \to \infty$: $x \to 0$.
        \item \textbf{Perspectiva bayesiana:} el término ridge equivale a un
          prior Gaussiano con media cero e igual varianza en todas las
          entradas (Cap.~7).
      \end{itemize}
    \end{column}
    \begin{column}{.48\textwidth}
      \begin{figure}
        \centering
        \includegraphics[width=\textwidth]{images/fig_4_13.png}
      \end{figure}
      \vspace{-2mm}
      {\scriptsize\color{gray}Fig.~4.13: menos varianza (elipses oscuras
      dentro de las claras), pero la media se desplaza.}
    \end{column}
  \end{columns}
  \fuente
\end{frame}

\begin{frame}{Ridge sin invertir: la versión SVD}
  \small
  Sustituyendo $A = U\Sigma V^T$ en (4.39):

  \begin{align*}
    (A^TA+\alpha I)^{-1}
      &= (V\Sigma^2V^T + \alpha VIV^T)^{-1} \\
      &= \big(V(\Sigma^2+\alpha I)V^T\big)^{-1} \\
      &= V(\Sigma^2+\alpha I)^{-1}V^T .
  \end{align*}

  \begin{block}{Lectura}
    \small
    \textbf{Basta sumarle $\alpha$ a cada valor singular al cuadrado antes de
    invertir.} Nada de matrices $n\times n$: es un filtro sobre los $\sigma_j$,
    \[
      \frac{\sigma_j}{\sigma_j^2 + \alpha},
    \]
    que apaga suavemente las direcciones con $\sigma_j$ pequeño. Es la misma
    familia de \emph{filtros espectrales} que usa el truncamiento duro de la
    SVD, pero continuo en vez de abrupto.
  \end{block}
  \fuente
\end{frame}

\begin{frame}{LASSO: regularización $\ell_1$}
  \small
  \[
    x = \argmin_{x'} \; \|Ax'-b\|_2^2 + \lambda\|x'\|_1 \tag{4.40}
  \]
  \begin{columns}[T,onlytextwidth]
    \begin{column}{.52\textwidth}
      \begin{itemize}
        \item \emph{Least Absolute Shrinkage and Selection Operator}
          (Tibshirani, 1996).
        \item Ridge \textbf{encoge} los coeficientes; LASSO los
          \textbf{apaga}: devuelve un $x$ \textbf{disperso}.
        \item $\lambda$ se elige por \textbf{validación cruzada}
          (típicamente 80/20).
        \item Motivación: interpretabilidad. En modelos de enfermedad hay
          decenas de predictores y la audiencia no es técnica.
        \item Antecedentes: soft-thresholding de Donoho--Johnstone; garrote de
          Breiman; navaja de Occam.
      \end{itemize}
    \end{column}
    \begin{column}{.46\textwidth}
      \begin{figure}
        \centering
        \includegraphics[width=\textwidth]{images/fig_4_14.png}
      \end{figure}
      \vspace{-2mm}
      {\scriptsize\color{gray}Fig.~4.14: los vértices de la bola $\ell_1$
      atraen la solución hacia los ejes.}
    \end{column}
  \end{columns}
  \fuente
\end{frame}

\begin{frame}{Elastic net y el mapa de las penalizaciones}
  \small
  \begin{columns}[T,onlytextwidth]
    \begin{column}{.52\textwidth}
      \textbf{Elastic net} (Zou--Hastie, 2005) combina ambas:
      \[
        x = \argmin_{x'}\; \|Ax'-b\|_2^2 + \lambda\|x'\|_1
              + \alpha\|x'\|_2^2 \tag{4.41}
      \]
      \begin{itemize}
        \item Ridge solo: encoge, no selecciona.
        \item LASSO solo: selecciona, pero con predictores muy
          correlacionados elige uno al azar y descarta el resto.
        \item Elastic net: selecciona y además \textbf{agrupa} predictores
          correlacionados.
      \end{itemize}
    \end{column}
    \begin{column}{.46\textwidth}
      \begin{figure}
        \centering
        \includegraphics[width=0.78\textwidth]{images/fig_4_15.png}
      \end{figure}
      \vspace{-2mm}
      {\scriptsize\color{gray}Fig.~4.15: conjuntos de nivel del elastic net con
      $\lambda=\alpha=1$; se parecen a la bola $\ell_{3/2}$.}
    \end{column}
  \end{columns}

  \vspace{1mm}
  \footnotesize
  Sin solución cerrada: se resuelven con métodos proximales, coordinate descent
  o solvers de QP/SOCP.
  \fuente
\end{frame}

% =============================================================================
\section{Restricciones y programación cuadrática}
% =============================================================================

\begin{frame}{Mínimos cuadrados con restricciones de igualdad}
  \small
  \[
    \min_x \; \|Ax-b\|_2^2
    \qquad \text{s.a. } Cx = d \tag{4.42}
  \]
  Equivalentemente, como \textbf{programa cuadrático}:
  \[
    \min_x \; \tfrac12 x^TQx + c^Tx
    \quad\text{s.a. } Cx=d,
    \qquad Q = 2A^TA,\;\; c^T = -2b^TA \tag{4.43--4.44}
  \]

  \textbf{Lagrangiano aumentado} (4.45):
  \[
    L(x,\lambda) = \tfrac12 x^TQx + c^Tx + \lambda^T(Cx-d)
  \]

  Condiciones de optimalidad $\Rightarrow$ sistema \textbf{KKT}:
  \[
    \begin{bmatrix} Q & C^T \\ C & 0 \end{bmatrix}
    \begin{bmatrix} x \\ \lambda \end{bmatrix}
    =
    \begin{bmatrix} -c \\ d \end{bmatrix}
    \qquad\Longleftrightarrow\qquad
    \begin{bmatrix} 2A^TA & C^T \\ C & 0 \end{bmatrix}
    \begin{bmatrix} x \\ \lambda \end{bmatrix}
    =
    \begin{bmatrix} 2A^Tb \\ d \end{bmatrix} \tag{4.47--4.48}
  \]
  \vspace{-2mm}
  {\footnotesize Se resuelve por SVD o QR. Sin restricciones
  ($C=0$) se recupera el problema estándar.}
  \fuente
\end{frame}

\begin{frame}{La norma mínima como caso particular}
  \small
  Con $A=I$ y $b=0$, los mínimos cuadrados restringidos se reducen a
  \textbf{norma mínima}:
  \[
    \min_x \|x\|_2^2 \quad\text{s.a. } Cx=d
    \qquad\Longrightarrow\qquad
    \begin{bmatrix} I & C^T \\ C & 0 \end{bmatrix}
    \begin{bmatrix} x \\ \lambda \end{bmatrix}
    =
    \begin{bmatrix} 0 \\ d \end{bmatrix} \tag{4.50}
  \]
  \[
    x^* = C^T(CC^T)^{-1}d
  \]

  \vspace{1mm}
  \begin{itemize}
    \item Todo el capítulo es un mismo sistema KKT con distintas matrices:
      cambio de $Q$, de $C$ o de $d$.
    \item \textbf{Sistemas con punto de silla:} la matriz KKT es simétrica pero
      \textbf{indefinida}, aunque $Q$ sea definida positiva. Los solvers
      genéricos pueden necesitar formulaciones especiales.
    \item Restricciones de \textbf{desigualdad} ($Cx \le d$, cajas, no
      negatividad) ya no caben en este esquema lineal: requieren QP ($\S4{.}6$)
      o el capítulo 6.
  \end{itemize}
  \fuente
\end{frame}

\begin{frame}{Qué es un programa cuadrático}
  \small
  \[
    \min_x \; \tfrac12 x^TQx + c^Tx
    \qquad \text{s.a. } Ax \le b \tag{4.51}
  \]
  \begin{columns}[T,onlytextwidth]
    \begin{column}{.50\textwidth}
      \begin{itemize}
        \item $Q$ simétrica.
        \item $Q \succeq 0$: QP \textbf{convexo}; todo mínimo local es global.
        \item $Q$ indefinida: no convexo, mucho más difícil.
        \item A diferencia del PL, el óptimo \textbf{no está en un vértice}:
          puede caer en una arista, una cara o el interior.
      \end{itemize}
    \end{column}
    \begin{column}{.48\textwidth}
      \begin{figure}
        \centering
        \includegraphics[width=\textwidth]{images/fig_4_16.png}
      \end{figure}
      \vspace{-2mm}
      {\scriptsize\color{gray}Fig.~4.16: restricciones inactivas (izq.: óptimo
      interior) y activas (der.: óptimo en la frontera).}
    \end{column}
  \end{columns}
  \fuente
\end{frame}

\begin{frame}{Tres familias de algoritmos}
  \small
  \begin{block}{1. Conjunto activo (\emph{active set})}
    Adivinar qué restricciones están activas, resolver el subproblema de
    igualdades con KKT y corregir el conjunto. \textbf{El simplex de Dantzig es
    un método de conjunto activo.}
  \end{block}
  \begin{block}{2. Punto interior (IPM)}
    Siguen una trayectoria central desde el interior del conjunto factible.
    \textbf{Escalan bien} a alta dimensión con restricciones lineales; son el
    estándar industrial para PL y QP grandes.
  \end{block}
  \begin{block}{3. Programación cuadrática secuencial (SQP)}
    Para problemas \textbf{no lineales}: en cada iteración se aproxima la
    curvatura con el Hessiano y las restricciones por linealización, y se
    resuelve un QP. Convergencia \textbf{superlineal} cerca de la solución.
  \end{block}
  \vspace{1mm}
  {\footnotesize El QP es el ladrillo con el que se construyen los métodos
  modernos para problemas no lineales restringidos.}
  \fuente
\end{frame}

\begin{frame}{Mínimos cuadrados \emph{son} un QP}
  \small
  \[
    \|Ax-b\|_2^2 = (Ax-b)^T(Ax-b)
                 = x^TA^TAx - 2b^TAx + b^Tb \tag{4.52}
  \]
  \[
    Q = A^TA, \qquad c^T = -2b^TA \tag{4.53}
  \]

  \begin{itemize}
    \item $A^TA$ es \textbf{simétrica y semidefinida positiva} para cualquier
      $A$: el QP resultante es \textbf{siempre convexo}.
    \item Cualquier restricción lineal (cajas, no negatividad, suma fija)
      se añade directamente y el problema sigue siendo un QP convexo.
  \end{itemize}

  \begin{block}{En la práctica: \texttt{cvxpy}}
    \small
    Reconoce automáticamente el QP y elige un solver especializado (OSQP,
    SCS, ECOS, GUROBI). Para QP grandes conviene \emph{warm starts}:
    \texttt{prob.solve(solver=cp.OSQP, warm\_start=True)}.
  \end{block}
  \fuente
\end{frame}

\begin{frame}[fragile]{Ejemplo: portafolio de Markowitz (E4.5)}
  \small
  \[
    \min_x \; \tfrac12 x^T\Sigma x
    \quad\text{s.a.}\quad
    r^Tx \ge R_0,\;\; \mathbf{1}^Tx = 1,\;\; x \ge 0
  \]
  \begin{columns}[T,onlytextwidth]
    \begin{column}{.52\textwidth}
      \begin{semiverbatim}
x = cp.Variable(n, nonneg=True)
obj = cp.Minimize(0.5*cp.quad_form(x,
                                 Sigma))
cons = [cp.sum(x) == 1,
        r @ x >= R0]
prob = cp.Problem(obj, cons)
prob.solve(solver=cp.SCS)
      \end{semiverbatim}
    \end{column}
    \begin{column}{.44\textwidth}
      \footnotesize
      \begin{itemize}
        \item $\Sigma$: covarianza de los rendimientos (el riesgo).
        \item $r$: rendimientos esperados.
        \item $x\ge 0$: prohibición de \textbf{ventas en corto}.
        \item Es \emph{el} QP canónico de las finanzas.
      \end{itemize}
    \end{column}
  \end{columns}
  \fuente
\end{frame}

% =============================================================================
\section{Mínimos cuadrados no lineales}
% =============================================================================

\begin{frame}{Gauss--Newton: el problema}
  \small
  \[
    L(\theta) = \tfrac12\sum_{j=1}^{m}\|y_j - N(x_j;\theta)\|_2^2
              = \tfrac12\sum_{j=1}^{m} f_j(\theta)^2
              = \tfrac12\|\mathbf{f}(\theta)\|_2^2 \tag{4.54}
  \]
  \begin{itemize}
    \item $\mathbf{f}(\theta)$: \textbf{vector de residuos}; $f_j$ es el error
      residual en el dato $j$.
    \item $\theta$: parámetros del modelo no lineal.
    \item Sin solución cerrada y \textbf{no convexo}.
    \item Es la formulación que cubre el \textbf{entrenamiento de redes
      neuronales}: $N$ es la red y $\theta$ son pesos y sesgos
      (Cap.~10).
    \item Se ataca con Newton, pero el Hessiano exacto es caro.
  \end{itemize}

  \vspace{1mm}
  \begin{block}{Estrategia}
    \scriptsize
    Resolver $\nabla_\theta L = 0$ requiere un sistema no lineal de ecuaciones.
    \textbf{Gauss--Newton aproxima el Hessiano} usando solo primeras derivadas,
    de modo que cada iteración es un problema de mínimos cuadrados
    lineales.
  \end{block}
  \fuente
\end{frame}

\begin{frame}{Gradiente, Jacobiano y Hessiano}
  \small
  \[
    \nabla_\theta L = \sum_{j=1}^{m} f_j(\theta)\nabla f_j(\theta)
                    = J^T\mathbf{f} \tag{4.55}
  \]
  \[
    J = D\mathbf{f} =
    \begin{bmatrix}
      \nabla f_1^T \\ \nabla f_2^T \\ \vdots \\ \nabla f_m^T
    \end{bmatrix}
    \in \R^{m\times n} \tag{4.56}
  \]
  \[
    H = \nabla_\theta^2 L
      = J^TJ + \sum_{j=1}^{m} f_j(\theta)\nabla^2 f_j(\theta)
      = J^TJ + Q \tag{4.58}
  \]

  \begin{block}{La suposición central de Gauss--Newton}
    \small
    Si los residuos $f_j$ son pequeños, el término $Q$ (segundas derivadas) es
    despreciable:
    \[
      H \approx J^TJ \tag{4.59}
    \]
    $J^TJ$ es simétrica y \textbf{semidefinida positiva}: garantiza una
    dirección de descenso.
  \end{block}
  \fuente
\end{frame}

\begin{frame}{La iteración de Gauss--Newton}
  \small
  Newton sería $\theta^{k+1} = \theta^k - H^{-1}\nabla L$, pero nunca se forma
  $H^{-1}$. Se resuelve el sistema y se da el paso:

  \[
    J_k^TJ_k\, s^k = -J_k^T\mathbf{f}(\theta^k) \qquad\text{\scriptsize\textcolor{gray}{(4.64a)}}
    \qquad\qquad
    \theta^{k+1} = \theta^k + s^k \qquad\text{\scriptsize\textcolor{gray}{(4.64b)}}
  \]

  \begin{block}{Equivalencia con un mínimos cuadrados}
    \small
    Resolver (4.64a) equivale exactamente a
    \[
      \min_s \; \tfrac12\|Js + \mathbf{f}\|_2^2 \tag{4.66}
    \]
    \textbf{Cada iteración de Gauss--Newton es un problema de mínimos cuadrados
    lineales} --- el del capítulo, con $A=J$ y $b=-\mathbf{f}$.
  \end{block}

  \vspace{1mm}
  \footnotesize
  Consecuencia: todo lo anterior se reutiliza dentro de Gauss--Newton (SVD,
  QR, regularización, gradiente conjugado).
  \fuente
\end{frame}

\begin{frame}{Los modos de falla}
  \small
  \begin{alertblock}{Cuándo Gauss--Newton se rompe}
    \begin{itemize}
      \item \textbf{Residuos grandes:} la aproximación $H \approx J^TJ$ ignora
        términos de segundo orden que ya no son pequeños
        $\Rightarrow$ \textbf{divergencia lejos de la solución}.
      \item \textbf{$J^TJ$ singular o casi singular:}
        \begin{enumerate}
          \item el paso $s^k$ puede ser enorme y \emph{pasarse} de la solución
            (overshoot);
          \item el sistema lineal está mal condicionado: pequeños errores en
            $J$ o $\mathbf{f}$ producen pasos muy distintos;
          \item la dirección puede no ser de descenso y la iteración se
            \emph{estanca}.
        \end{enumerate}
      \item Salvo residuos muy pequeños, la convergencia es solo
        \textbf{de primer orden}, no cuadrática.
    \end{itemize}
  \end{alertblock}

  \vspace{1mm}
  Vasta literatura dedicada a estas patologías. La cura práctica estándar es
  \textbf{amortiguar} el Hessiano.
  \fuente
\end{frame}

\begin{frame}{Levenberg--Marquardt: amortiguar el Hessiano}
  \small
  \[
    H \approx J^TJ + \alpha I \qquad\text{\scriptsize\textcolor{gray}{(4.73)}}
    \qquad\Longrightarrow\qquad
    (J^TJ + \alpha I)\,\Delta\theta = -J^T\mathbf{f} \qquad\text{\scriptsize\textcolor{gray}{(4.74)}}
  \]

  \begin{columns}[T,onlytextwidth]
    \begin{column}{.50\textwidth}
      \begin{itemize}
        \item $\alpha = 0$: \textbf{Gauss--Newton} puro (rápido cerca del
          óptimo).
        \item $\alpha \to \infty$: $\Delta\theta = -\frac{1}{\alpha}J^T\mathbf{f}$:
          \textbf{descenso de gradiente} con paso $1/\alpha$ (lento pero
          seguro).
        \item $\alpha$ \textbf{interpola} entre exploración y explotación.
      \end{itemize}
    \end{column}
    \begin{column}{.48\textwidth}
      \begin{block}{Dos lecturas}
        \scriptsize
        \textbf{1.} Es \textbf{regularización de Tikhonov} ($\S4{.}4$) aplicada
        al paso de Newton: suma $\alpha$ a los valores singulares al cuadrado.

        \vspace{2mm}
        \textbf{2.} Es una \textbf{región de confianza}: $\alpha$ controla
        cuánto se confía en el modelo cuadrático local.
      \end{block}
    \end{column}
  \end{columns}

  \vspace{1mm}
  \footnotesize
  Usado en visión por computadora, ajuste de curvas, robótica y estimación de
  parámetros en control. Es el ``mínimos cuadrados no lineales'' por defecto.
  \fuente
\end{frame}

\begin{frame}{Cómo se adapta $\alpha$}
  \small
  \begin{itemize}
    \item Comparar la reducción \textbf{real} de la pérdida con la
      \textbf{predicha} por el modelo cuadrático:
      \[
        \rho = \frac{\Delta L_{\text{actual}}}{\Delta L_{\text{predicho}}}
        \tag{4.76}
      \]
      \[
        \Delta L_{\text{actual}} = L(\theta^k)-L(\theta^{k+1}) \tag{4.77}
      \]
      \[
        \Delta L_{\text{predicho}}
        = -J^T\mathbf{f}\,\Delta\theta
          - \Delta\theta^T(J^TJ+\alpha I)\Delta\theta \tag{4.78}
      \]
    \item Regla típica: $\rho > 0.75$ $\Rightarrow$ reducir $\alpha$;
      $\rho < 0.25$ $\Rightarrow$ aumentar $\alpha$.
      Multiplicar o dividir por un factor $\nu > 1$.
    \item Se usa el cociente, no la magnitud absoluta: cerca del mínimo las
      reducciones son diminutas pero deben seguir siendo las predichas.
    \item Inicialización habitual: $\alpha \sim 10^{-3}$.
    \item Variante más fina: escalar cada entrada diagonal,
      $H = J^TJ + \alpha\,\mathrm{diag}(J^TJ)$ (4.80), para respetar la
      \textbf{curvatura local} de cada parámetro.
  \end{itemize}
  \fuente
\end{frame}

% =============================================================================
\section{El método del gradiente conjugado}
% =============================================================================

\begin{frame}{Gradiente conjugado: para qué}
  \small
  \[
    Ax = b, \qquad A \in \R^{n\times n} \text{ simétrica definida positiva}
    \tag{4.81}
  \]
  \begin{itemize}
    \item \textbf{No factoriza $A$}: solo necesita productos matriz--vector.
    \item Ideal para sistemas \textbf{grandes y dispersos}; el ejemplo clásico es
      la ecuación de Poisson en dinámica de fluidos.
    \item Convergencia garantizada en \textbf{a lo sumo $n$ iteraciones}.
    \item Su uso estrella aquí: resolver la ecuación de Gauss--Newton
      \[
        J^TJ\,\Delta x = -J^T\mathbf{f} \tag{4.82}
      \]
      \textbf{sin formar ni almacenar $J^TJ$}. La memoria se dispara:
      $J$ es $m\times n$ y $J^TJ$ sería $n\times n$ densa.
    \item También mejora la estabilidad frente a la inversión explícita de un
      Hessiano mal condicionado. El precio: la convergencia sigue dependiendo
      de $\kappa(J^TJ)$.
  \end{itemize}
  \fuente
\end{frame}

\begin{frame}{La idea: base $A$-conjugada}
  \small
  Direcciones \textbf{conjugadas} = ortogonales en el producto interior
  $\langle u,v\rangle_A = u^TAv$:
  \[
    p_i^TAp_j = 0 \qquad \forall\, i \ne j \tag{4.83}
  \]
  Si $A \succ 0$, el conjunto $\{p_j\}_{j=1}^n$ es una \textbf{base de
  $\R^n$}. Se expande la solución y se despeja cada coeficiente:
  \[
    x = \sum_{j=1}^{n}\alpha_j p_j \tag{4.84}
  \]
  multiplicando $Ax=b$ por $p_k^T$ y usando la conjugación, todos los términos
  salvo uno se anulan:
  \[
    \alpha_k = \frac{p_k^Tb}{p_k^TAp_k} \tag{4.88}
  \]
  \begin{block}{Por qué es mejor que el descenso de gradiente}
    \scriptsize
    En CG, minimizar a lo largo de una dirección \textbf{no destruye} el progreso
    hecho en las anteriores. El descenso de gradiente sí lo destruye: por eso
    zigzaguea.
  \end{block}
  \fuente
\end{frame}

\begin{frame}{El algoritmo}
  \small
  Inicializar e iterar:
  \[
    x_0,\qquad r_0 = b-Ax_0,\qquad p_0 = r_0 \tag{4.94}
  \]
  \[
    \alpha_k = \frac{r_k^Tr_k}{p_k^TAp_k} \qquad\text{\scriptsize\textcolor{gray}{(4.95a)}}
    \qquad
    x^{k+1} = x^k + \alpha_k p_k \qquad\text{\scriptsize\textcolor{gray}{(4.95b)}}
    \qquad
    r_{k+1} = r_k - \alpha_k Ap_k \qquad\text{\scriptsize\textcolor{gray}{(4.95c)}}
  \]
  \[
    \beta_k = \frac{r_{k+1}^Tr_{k+1}}{r_k^Tr_k} \qquad\text{\scriptsize\textcolor{gray}{(4.96a)}}
    \qquad
    p_{k+1} = r_{k+1} + \beta_k p_k \qquad\text{\scriptsize\textcolor{gray}{(4.96b)}}
  \]

  \begin{columns}[T,onlytextwidth]
    \begin{column}{.48\textwidth}
      \textbf{Qué guarda:} solo tres vectores ($x^k$, $r_k$, $p_k$).
      Memoria $\Theta(n)$.
    \end{column}
    \begin{column}{.48\textwidth}
      \textbf{Qué cuesta por iteración:} un producto $Ap_k$, dos productos
      internos y tres operaciones vectoriales. Nada de factorizar.
    \end{column}
  \end{columns}
  \fuente
\end{frame}

\begin{frame}{Convergencia, precondicionamiento y GMRES}
  \small
  \[
    \|x^k - x^*\|_A \le C\left(\frac{\kappa(A)-1}{\kappa(A)+1}\right)^{k}
    \tag{4.97}
  \]
  \begin{itemize}
    \item Cota en la \textbf{norma de energía} $\|\cdot\|_A$; $\kappa$ es el
      cociente de valores propios extremos.
    \item Con $\kappa$ grande la cota es casi 1 y CG se arrastra:
      \textbf{el precondicionamiento es obligatorio en la práctica}.
    \item Se resuelve $M^{-1}Ax = M^{-1}b$ (4.98) con $M \approx A$ pero fácil
      de invertir.
  \end{itemize}

  \begin{block}{Panel 4.4: GMRES}
    \scriptsize
    Cuando $A$ \textbf{no} es simétrica definida positiva, CG no aplica.
    \textbf{GMRES} minimiza $\|r\|_2$ sobre el subespacio de Krylov usando la
    iteración de Arnoldi. El costo de ortogonalización crece con las
    iteraciones, así que se usa reiniciado (\textbf{GMRES(m)}). También depende
    críticamente del precondicionador. Es el \emph{caballo de batalla} para
    sistemas no simétricos (fluidos, electromagnetismo).
  \end{block}
  \fuente
\end{frame}

% =============================================================================
\section{Cierre}
% =============================================================================

\begin{frame}{Puntos para recordar}
  \small
  \begin{enumerate}
    \item \textbf{Dos regímenes, una matriz:} sobredeterminado $\to$ mínimos
      cuadrados; subdeterminado $\to$ norma mínima. Ambas son $A^{\dagger}b$.
    \item \textbf{Nunca $A^TA$:} se usa la SVD. Es estable, siempre existe y da
      la aproximación óptima de bajo rango (Eckart--Young).
    \item \textbf{Conditioning es el enemigo:} el ruido se multiplica por
      $\kappa(A)$. Truncar la SVD o regularizar (ridge) son las dos curas.
    \item \textbf{$\ell_1$ frente a $\ell_2$:} $\ell_1$ da robustez a outliers
      (mediana) y dispersidad (LASSO, compressed sensing). El precio es perder
      la fórmula cerrada.
    \item \textbf{Las restricciones son un sistema KKT:} una sola estructura
      matricial cubre norma mínima y mínimos cuadrados restringidos.
    \item \textbf{Mínimos cuadrados es un QP convexo:} cualquier restricción
      lineal encaja; \texttt{cvxpy} lo resuelve.
    \item \textbf{No lineal:} Gauss--Newton convierte cada iteración en un
      mínimos cuadrados; Levenberg--Marquardt lo amortigua para no divergir.
    \item \textbf{Grande y disperso:} gradiente conjugado, con productos
      matriz--vector y memoria $\Theta(n)$.
  \end{enumerate}
  \fuente
\end{frame}

\begin{frame}{El mapa completo}
  \centering
  \small
  \vspace{2mm}
  \begin{tabular}{@{}p{0.30\textwidth}p{0.62\textwidth}@{}}
    \toprule
    \textbf{Herramienta} & \textbf{Qué problema resuelve} \\
    \midrule
    Ecuaciones normales & LS sobredeterminado, rango completo, bien condicionado \\
    SVD / pseudo-inversa & LS y norma mínima, estable, cualquier rango \\
    Truncamiento SVD & conditioning malo, ruido en los modos débiles \\
    Ridge (Tikhonov) & conditioning malo con sesgo controlado \\
    LASSO / elastic net & selección de variables, modelos interpretables \\
    KKT & restricciones de igualdad \\
    QP (active set / IPM) & restricciones de desigualdad, cajas, portafolios \\
    Gauss--Newton & regresión no lineal, redes neuronales \\
    Levenberg--Marquardt & Gauss--Newton con residuos grandes o mal condicionado \\
    Gradiente conjugado & sistemas enormes y dispersos, sin factorizar \\
    \bottomrule
  \end{tabular}

  \vspace{4mm}
  {\footnotesize\color{gray}
    Todo el capítulo es una sola idea con distintos disfraces: proyectar
    $b$ sobre un subespacio y controlar cuánto se confía en las direcciones
    débiles.}
\end{frame}

\begin{frame}{Para profundizar}
  \small
  \textbf{Libros}
  \begin{itemize}
    \item Brunton \& Kutz, \emph{Data-Driven Science and Engineering} (2022).
    \item James, Witten, Hastie \& Tibshirani, \emph{An Introduction to
      Statistical Learning} (2013).
    \item Trefethen, \emph{Numerical Linear Algebra} (2022).
    \item Antoulas, \emph{Approximation of Large-Scale Dynamical Systems}
      (2005).
  \end{itemize}

  \textbf{Artículos fundacionales}
  \begin{itemize}
    \item Tikhonov (1943), \emph{On the stability of inverse problems}.
    \item Tibshirani (1996), \emph{Regression shrinkage and selection via the
      Lasso}.
    \item Zou \& Hastie (2005), \emph{Regularization and variable selection via
      the elastic net}.
    \item Donoho (2006), \emph{Compressed sensing}.
    \item Cand\`es, Romberg \& Tao (2006), \emph{Stable signal recovery from
      incomplete and inaccurate measurements}.
  \end{itemize}

  \vfill
  \centering
  {\footnotesize\color{gray}
    Siguiente capítulo: optimización no suave y global.}
\end{frame}

% =============================================================================
\end{document}
